\section{Future Work}
\label{sec:future}

The most important next steps are output correctness validation and extended benchmarking.

\textbf{VAF accuracy validation (done).}
A chi-squared goodness-of-fit test comparing observed alt-read counts against $\text{Binomial}(D, \text{VAF})$ now validates the Bernoulli sampling model at four VAF levels, with a companion power test confirming rejection of deterministic spike-in.
The next step is end-to-end validation: generate reads at a known VAF, align with BWA, count alt reads in pileup, and compare the observed VAF distribution against the configured value.

\textbf{cfDNA fragment distribution validation.}
The cfDNA mode uses a three-component Gaussian mixture.
Extracting insert sizes from the BAM output and comparing to the configured parameters would confirm that the fragment model is applied correctly.

\textbf{Downstream variant caller benchmarking.}
The downstream variant calling benchmark described in Section~\ref{sec:vc} validates VarForge data against a real four-stage pipeline (bwa-mem2, Mutect2, rtg-tools).
Extending this to additional callers (VarDict, Strelka2) and UMI-aware deduplication pipelines (fgbio, HUMID) would provide a more complete sensitivity-precision landscape across tools.

\textbf{Compute-bound thread scaling.}
Current thread scaling data covers only the I/O-bound 10\,MB, 30x workload where threading degrades performance.
Testing thread scaling for UMI simulation (where per-pair compute cost is higher) would characterise when Rayon parallelism provides real benefit.

\textbf{Long-read sequencing.}
PacBio HiFi and Nanopore R10 presets now include sequencing error model parameters, and \texttt{learn-profile} can extract per-cycle error rates and indel rates from long-read BAMs.
The remaining gap is long-read-specific error structure: homopolymer run-length errors (Nanopore) and CCS-specific artefacts (PacBio) are not yet modelled explicitly.
Dedicated long-read error passes would improve realism for long-read variant caller benchmarking.

\textbf{Enhanced error profiles (partially addressed).}
The v0.2.0 \texttt{ErrorOrchestrator} adds cycle-position errors, k-mer context errors, sequencing indels, strand bias, and phasing bursts.
The \texttt{learn-profile} subcommand now extracts per-cycle error rates and indel rates from BAMs, enabling data-driven parameterisation.
The remaining gap is the full joint correlation structure modelled by ReSeq\autocite{schmeing2021reseq}: pairwise quality dependencies across read positions are not yet captured.
